\appendixsection{一个连通性性质}

引理 1. 设 X 为连通拓扑空间，\(\Omega\) 为 X 的稠密开子集。若对任意 \(x \in X\) ，存在 x 的邻域 V 使得 \(V \cap \Omega\) 连通，则 \(\Omega\) 连通。

事实上，设 \(\Omega_0\) 是 \(\Omega\) 的非空既开又闭的子集。取 \(x \in X\) ，并设 \(V\) 是 \(x\) 的邻域，使得 \(V \cap \Omega\) 连通。若 \(x \in \overline{\Omega}_0\) ，

\[
(\mathrm{V} \cap \Omega) \cap \Omega_ {0} = \mathrm{V} \cap \Omega_ {0} \neq \varnothing ,
\]

于是 \(V \cap \Omega \subset \Omega_{0}\) 。这样，由于 \(\Omega\) 在 X 中稠密，\(\overline{\Omega}_{0}\) 是 x 的邻域。因而 \(\overline{\Omega}_{0}\) 非空、既开又闭；又因 X 连通，故 \(\overline{\Omega}_{0} = X\) 。由于 \(\Omega_{0}\) 在 \(\Omega\) 中闭，由此推出 \(\Omega_{0} = \Omega \cap \overline{\Omega}_{0} = \Omega\) ，这就证明了 \(\Omega\) 连通。

引理 2. 设 U 为 \(\mathbf{C}^n\) 中的开球，\(f: \mathrm{U} \to \mathbf{C}\) 为不恒为零的全纯函数。设 A 为 U 的子集，使 \(f = 0\) 在 A 上成立。则 U-A 在 U 中稠密且连通。

U-A 的稠密性由《微分与解析流形（结果）》3.2.5 得出。先设 \(n = 1\) 。若 \(a \in \mathrm{A}\) ，则 \(f\) 在 \(a\) 处的幂级数展开（《微分与解析流形（结果）》3.2.1）不恒为 0，由此推出，在 U 中存在邻域 \(\mathrm{V}_a\) 包含 \(a\) ，使 \(f\) 在 \(\mathrm{V}_a - \{a\}\) 上不取零值。于是 \(a\) 在 A 中孤立，这就证明 A 是 U 的离散子集，从而由于 U 在无穷远处可数，A 是可数的。取 \(x, y \in \mathrm{U} - \mathrm{A}\) 。把 \(x\) （相应地 \(y\) ）与 A 中一点连接起来的实仿射直线的并集是贫集（《一般拓扑学》第九章 §5 第 2 节 定义 2）。因此存在 \(z \in \mathrm{U} - \mathrm{A}\) ，使线段 \([x, z]\) 与 \([y, z]\) 都不与 A 相交。于是三点 \(x, y, z\) 属于 U-A 的同一连通分支，这就证明了 \(n = 1\) 情形的引理。转而考虑一般情形。可以假设 A 是 \(f\) 的零点集（《一般拓扑学》第一章 §11 第 1 节 命题 1）。取 \(x, y \in \mathrm{U} - \mathrm{A}\) ，并设 L 为包含 \(x\) 与 \(y\) 的仿射直线。\(f\) 在 \(\mathrm{L} \cap \mathrm{U}\) 上的限制不恒为零，因为 \(x \in \mathrm{L} \cap \mathrm{U}\) 。由已证者，\(x\) 与 \(y\) 属于 \((\mathrm{L} \cap \mathrm{U}) - (\mathrm{L} \cap \mathrm{A})\) 的同一连通分支，从而属于 U-A 的同一连通分支。

引理 3. 设 X 为有限维连通复解析流形，并设 A 为 X 的满足下述条件的子集：

对任意 \(x \in X\) ，存在在 x 处不取零值的解析函数芽 \(f_{x}\) ，使得 A 在 x 处的芽包含于 \(f_{x}\) 的零点集在 x 处的芽之中。

则 X-A 在 X 中稠密且连通。

X-A 的稠密性由《微分与解析流形（结果）》3.2.5 得出。可以假设 A 是闭的（《一般拓扑学》第一章 §11 第 1 节 命题 1）。对任意 \(x \in \mathrm{X}\) ，存在开邻域 \(\mathrm{V}\) 包含 \(x\) ，以及同构 \(c\) ，它把 \(\mathrm{V}\) 映到 \(\mathbf{C}^n\) 中的一个开球，且使得 \(c(\mathrm{A} \cap \mathrm{V})\) 包含于某个在 \(c(\mathrm{V})\) 上不恒为零的全纯函数的零点集中。于是由引理 2，\(\mathrm{V} \cap (\mathrm{X} - \mathrm{A})\) 连通。依据引理 1，这就证明 \(\mathrm{X} - \mathrm{A}\) 连通。

